% ===============================================================
\section{Introduction and organisation of the work}
% ===============================================================

This report documents the work done for Project~2 of ME249. In Part~1, a genetic algorithm (GA) is used to
extract, from nucleate boiling data obtained for a water/2-propanol mixture at gravity levels between
$0.01\,\gen$ and $2\,\gen$, the dependence of the heat flux on wall superheat, gravity, Marangoni
(surface-tension) effects and pressure. In Part~2, the same GA is adapted to calibrate a simple
performance model of a heat-pipe heat exchanger that cools an electronics cabinet.

\subsection{Division of the work}
\label{sec:division}

\textbf{I worked alone on this project}, so all tasks (programming, runs, analysis, figures and report)
were done by me. The adjustment sheet for individuals working alone allows omitting part~(b) of
Task~2.1 and Task~2.3. The two extra tasks ET1 and ET2 are assigned to teams of three. \textbf{I still
completed the whole project}: Tasks~2.1(a) and (b), 2.2, 2.3, ET1, ET2, 2.4, 2.5 and 2.6.

\begin{table}[H]
\centering\small
\caption{Tasks, programs and where the results are reported. All programs are Jupyter notebooks in the
folder \code{code/}. Their full listings are in Appendix~\ref{app:code}.}
\label{tab:tasks}
\begin{tabularx}{\linewidth}{@{}l l X l@{}}
\toprule
Task & Program (notebook) & Content & Section\\
\midrule
2.1 (a), (b) & \code{P2\_Task2.1.ipynb}      & data check, log-log plot of $\qpp$ vs.\ superheat & \ref{sec:t21}\\
2.2          & \code{P2\_Task2.2.ipynb}      & three-constant model, Eq.~(1.1), 45 points & \ref{sec:t22}\\
2.3          & \code{P2\_Task2.3.ipynb}      & five-constant model, Eq.~(1.5), 77 + 15 points & \ref{sec:t23}\\
ET1          & \code{P2\_ET1.ipynb}          & fractional sensitivity of $\qpp$ to its inputs & \ref{sec:et1}\\
ET2          & \code{P2\_ET2.ipynb}          & inverse relation for the superheat, Eq.~(ET2.1) & \ref{sec:et2}\\
2.4, 2.5     & \code{P2\_Task2.4-2.5.ipynb}  & heat-pipe model, $(UA)_e$, $(UA)_c$, $\alpha$ & \ref{sec:part2}\\
2.6          & this report                   & summary, code changes, conclusions & \ref{sec:conclusions}, \ref{sec:codechanges}\\
\bottomrule
\end{tabularx}
\end{table}

\subsection{Software and reproducibility}
All programs were written in Python~3 (NumPy, Matplotlib) as Jupyter notebooks. They are built on the
skeleton programs \code{CodeP2.1F26}, \code{CodeP2.2F26} and \code{CodeP2.3F26}. Every run is
reproducible because both random number generators are seeded (see Section~\ref{sec:codechanges}). Each
notebook writes its figures to \code{report/figures/} and its numerical results to \code{code/results/*.json}.
The tables of this report are generated from these files.
The best-fit constants are always the ones found by the genetic algorithm. In addition, a deterministic
minimisation (SciPy Nelder--Mead, multi-start) of \emph{exactly the same} error function is used
\emph{only as an independent check}: it tells us how close the GA constants are to the true minimum of
the error function.

\subsection{The genetic algorithm of the skeleton program}
\label{sec:ga}
The evolution loop of \code{CodeP2.2F26} was kept unchanged in all the programs. Each organism (``DNA strand'')
$j$ is a vector $[n_0=-1,\,n_1,\dots,n_5]$ of model constants. The population has $N_S=N_D$ organisms,
and at every generation $k$ the loop does the following:
\begin{enumerate}[label=\arabic*.]
  \item \textbf{Fitness.} Organism $i$ is compared with data point $i$ only. Its error
        $f_{err,i}$ (e.g.\ Eq.~(1.3)) gives the absolute fractional error
        $\mathrm{aFerr}_i=|f_{err,i}|/|\ln\qpp_{data,i}|$. Because mating and selection constantly
        reorder the population, the pairing between organisms and data points is random.
  \item \textbf{Selection.} The cut-off is $c_{lim}=\mathrm{MFRAC}\times\mathrm{median}(\mathrm{aFerr})$,
        with $\mathrm{MFRAC}=0.5$. The \code{nkeep} organisms with $\mathrm{aFerr}<c_{lim}$ survive
        unchanged.
  \item \textbf{Mating and mutation.} $N_S-\code{nkeep}$ offspring are created. For each offspring, two
        parents are drawn at random among the first $\code{nkeep}+1$ organisms of the previous population.
        Each constant is copied from one of the two parents (probability $1/2$) and multiplied by
        $1+\mathrm{MUT}\cdot 2(0.5-U)$, with $U$ uniform in $[0,1]$. The mutation is therefore uniform in
        $\pm\mathrm{MUT}$ (MUT~$=0.085$ in the skeleton code, $\pm9\,\%$ in the statement).
  \item \textbf{Population average.} The averages $\bar n_j(k)$ of the constants over the population are
        computed. The error of these average constants over \emph{all} data points,
        $\Ferrmean(k)=F_{err}/N_D$ (\code{aFerrmeanavgn}), measures the quality of the fit. Its minimum
        over all generations, \code{aFerrmeanavgnMin}, and the corresponding constants
        (\code{n1min}, \dots) are the best fit. Time averages of $\bar n_j$ for $k>800$ are also computed.
\end{enumerate}
Two properties of this algorithm explain most of the behaviour reported below. First, the selection
threshold is \emph{relative}: it compares errors to their median. Second, the mutations are
\emph{multiplicative}: a constant moves by at most a factor $(1\pm\mathrm{MUT})$ per generation, and it
can never change sign.
